% Coverage: Chapter 8, Section 8.6, physical PDF pp. 378--384
% (printed pp. 355--361; p. 378 begins with the end of Section 8.5).
% Semantic range: Section 8.6 heading through the final sentence referring to the
% chapter Appendix. Inventory: equations (8.6.1)--(8.6.27), two ordinary numbered
% footnotes, no figures, tables, problems, references, or unresolved uncertainties.

\subsection{Feynman Rules for Spinor Electrodynamics}

We are now in a position to state the Feynman rules for calculating the
$S$-matrix in quantum electrodynamics. For definiteness, we will consider the
electrodynamics of a single species of spin $\frac12$ particles of charge
$q=-e$ and mass $m$. We will call these fermions electrons, but the same
formalism applies to muons and other such particles. The simplest gauge- and
Lorentz-invariant Lagrangian for this theory is%
\footnote{In Chapter 12 we will discuss reasons why more complicated terms are
excluded from the Lagrangian density.}
\begin{equation}
\mathcal L=-\frac14F_{\mu\nu}F^{\mu\nu}
-\overline\Psi\bigl(\gamma^\mu[\partial_\mu+ieA_\mu]+m\bigr)\Psi.
\tag{8.6.1}\label{eq:8.6.1}
\end{equation}
The electric current four-vector is then simply
\begin{equation}
J^\mu\equiv\frac{\partial\mathcal L}{\partial A_\mu}
=-ie\overline\Psi\gamma^\mu\Psi.
\tag{8.6.2}\label{eq:8.6.2}
\end{equation}
The interaction \eqref{eq:8.4.23} in the interaction picture is here
\begin{equation}
H_I(t)=+ie\int d^3x\,
\bigl(\overline\psi(\mathbf x,t)\gamma^\mu\psi(\mathbf x,t)\bigr)
a_\mu(\mathbf x,t)+V_{\mathrm{Coul}}(t).
\tag{8.6.3}\label{eq:8.6.3}
\end{equation}
(There is no $H_{I,\mathrm{matter}}$ here.) As we have seen, the Coulomb term
$V_{\mathrm{Coul}}(t)$ just serves to cancel a part of the photon propagator
that is non-covariant and local in time.

Following the general results of Section 6.3, we can state the momentum space
Feynman rules for the connected part of the $S$-matrix in this theory as
follows:

\medskip
\noindent\textbf{(i)} Draw all Feynman diagrams with up to some given number
of vertices. The diagrams consist of electron lines carrying arrows and photon
lines without arrows, with the lines joined at vertices, at each of which there
is one incoming and one outgoing electron line and one photon line. There is
one external line coming into the diagram from below or going upwards out of
the diagram for each particle in the initial or final states, respectively;
electrons are represented by external lines carrying arrows pointing upwards
into or out of the diagram, while positrons are represented by lines carrying
arrows pointing downwards into or out of the diagram. There are also as many
internal lines as are needed to give each vertex the required number of
attached lines. Each internal line is labelled with an off-mass-shell
four-momentum flowing in a definite direction along the line (taken
conventionally to flow along the direction of the arrow for electron lines).
Each external line is labelled with the momentum and spin $z$-component or
helicity of the electron or photon in the initial and final states.

\medskip
\noindent\textbf{(ii)} Associate factors with the components of the diagram
as follows:

\begin{center}
\textit{Vertices}
\end{center}

Label each vertex with a four-component Dirac index $\alpha$ at the electron
line with its arrow coming into the vertex, a Dirac index $\beta$ at the
electron line with its arrow going out of the vertex, and a spacetime index
$\mu$ at the photon line. For each such vertex, include a factor
\begin{equation}
(2\pi)^4e(\gamma^\mu)_{\beta\alpha}
\delta^4(k-k'+q),
\tag{8.6.4}\label{eq:8.6.4}
\end{equation}
where $k$ and $k'$ are the electron four-momenta entering and leaving the
vertex, and $q$ is the photon four-momentum entering the vertex (or minus the
photon momentum leaving the vertex).

\begin{center}
\textit{External lines:}
\end{center}

Label each external line with the three-momentum $\mathbf p$ and spin
$z$-component or helicity $\sigma$ of the particle in the initial or final
state. For each line for an electron in the final state running out of a
vertex carrying a Dirac label $\beta$ on this line, include a factor
\footnote{A matrix $\beta$ has been extracted from the interaction in \eqref{eq:8.6.4},
so that $\overline u$ and $\overline v$ appear instead of $u^\dagger$ and
$v^\dagger$.}
\begin{equation}
\frac{\overline u_\beta(\mathbf p,\sigma)}{(2\pi)^{3/2}}.
\tag{8.6.5}\label{eq:8.6.5}
\end{equation}
For each line for a positron in the final state running into a vertex carrying
a Dirac label $\alpha$ on this line, include a factor
\begin{equation}
\frac{v_\alpha(\mathbf p,\sigma)}{(2\pi)^{3/2}}.
\tag{8.6.6}\label{eq:8.6.6}
\end{equation}
For each line for an electron in the initial state running into a vertex
carrying a Dirac label $\alpha$ on this line, include a factor
\begin{equation}
\frac{u_\alpha(\mathbf p,\sigma)}{(2\pi)^{3/2}}.
\tag{8.6.7}\label{eq:8.6.7}
\end{equation}
For each line for a positron in the initial state running out of a vertex
carrying a Dirac label $\beta$ on this line, include a factor
\begin{equation}
\frac{\overline v_\beta(\mathbf p,\sigma)}{(2\pi)^{3/2}}.
\tag{8.6.8}\label{eq:8.6.8}
\end{equation}
The $u$s and $v$s are the four-component spinors discussed in Section 5.5.
For each line for a photon in the final state connected to a vertex carrying a
spacetime label $\mu$ on this line, include a factor
\begin{equation}
\frac{e^*_\mu(\mathbf p,\sigma)}{(2\pi)^{3/2}\sqrt{2p^0}}.
\tag{8.6.9}\label{eq:8.6.9}
\end{equation}
For each line for a photon in the initial state connected to a vertex carrying
a spacetime label $\mu$ on this line, include a factor
\begin{equation}
\frac{e_\mu(\mathbf p,\sigma)}{(2\pi)^{3/2}\sqrt{2p^0}}.
\tag{8.6.10}\label{eq:8.6.10}
\end{equation}
The $e_\mu$ are the photon polarization four-vectors described in the previous
section.

\begin{center}
\textit{Internal lines:}
\end{center}

For each internal electron line carrying a four-momentum $k$ and running from
a vertex carrying a Dirac label $\beta$ to another vertex carrying a Dirac
label $\alpha$, include a factor
\begin{equation}
\frac{-i}{(2\pi)^4}
\frac{[-i\sl{k}+m]_{\alpha\beta}}{k^2+m^2-i\epsilon}.
\tag{8.6.11}\label{eq:8.6.11}
\end{equation}
(We are here using the very convenient `Dirac slash' notation; for any
four-vector $v^\mu$, $\sl{v}$ denotes $\gamma_\mu v^\mu$.) For each
internal photon line carrying a four-momentum $q$ that runs between two
vertices carrying spacetime labels $\mu$ and $\nu$ include a factor
\begin{equation}
\frac{-i}{(2\pi)^4}\frac{\eta_{\mu\nu}}{q^2-i\epsilon}.
\tag{8.6.12}\label{eq:8.6.12}
\end{equation}

\medskip
\noindent\textbf{(iii)} Integrate the product of all these factors over the
four-momenta carried by the internal lines, and sum over all Dirac and
spacetime indices.

\medskip
\noindent\textbf{(iv)} Add up the results obtained in this way from each
Feynman diagram.

Additional combinatoric factors and fermionic signs may need to be included,
as described in parts (v) and (vi) of Section 6.1.

The difficulty of evaluating Feynman diagrams increases rapidly with the
number of internal lines and vertices, so it is important to have some idea of
what numerical factors tend to suppress the contributions of the more
complicated diagrams. We shall estimate these numerical factors including not
only the factors of the electronic charge $e$ associated with vertices, but
also the factors of 2 and $\pi$ from vertices, propagators, and momentum space
integrals.

Consider a connected Feynman diagram with $V$ vertices, $I$ internal lines,
$E$ external lines, and $L$ loops. These quantities are not independent, but
are subject to relations already used in Section 6.3:
\[
L=I-V+1,\qquad 2I+E=3V.
\]
There is a factor $e(2\pi)^4$ from each vertex, a factor $(2\pi)^{-4}$
from each internal line, and a four-dimensional momentum space integral for
each loop. The volume element in four-dimensional Euclidean space in terms of
a radius parameter $\kappa$ is $\pi^2\kappa^2d\kappa^2$, so each loop
contributes a factor $\pi^2$. Thus the diagram will contain a factor
\[
(2\pi)^{4V}e^V(2\pi)^{-4I}\pi^{2L}
=(2\pi)^4e^{E-2}\left(\frac{e^2}{16\pi^2}\right)^L.
\]
The number $E$ of external lines is fixed for a given process, so we see that
the expansion parameter that governs the suppression of Feynman graphs for
each additional loop is
\[
\frac{e^2}{16\pi^2}=\frac{\alpha}{4\pi}=5.81\cdot10^{-4}.
\]
Fortunately this is small enough that good accuracy can usually be obtained
from Feynman diagrams with at most a few loops.

\begin{center}
$*\quad *\quad *$
\end{center}

We must say a little more about the spin states of photons and electrons in
realistic experiments, where not every particle in the initial and final
states has a definite known helicity or spin $z$-component. This consideration
is especially important for photons, which in practice are often characterized
by a state of transverse or elliptical polarization rather than helicity. As
we saw in the previous section, for photons of helicity $\mathord\pm1$, the
polarization vectors are
\[
e^\mu(\mathbf p,\mathord\pm1)=R(\widehat{\mathbf p})
\begin{bmatrix}
0\\[1pt]1/\sqrt2\\[1pt]\mathord\pm i/\sqrt2\\[1pt]0
\end{bmatrix},
\]
where $R(\widehat{\mathbf p})$ is the standard rotation that takes the
$z$-axis to the $\mathbf p$ direction. These are not the only possible photon
states; in general, a photon state can be a linear combination of helicity
states $\ket{\mathbf p,\mathord\pm1}$
\begin{equation}
\alpha_+\ket{\mathbf p,+1}+\alpha_-\ket{\mathbf p,-1}
\tag{8.6.13}\label{eq:8.6.13}
\end{equation}
which is properly normalized if
\begin{equation}
\lvert\alpha_+\rvert^2+\lvert\alpha_-\rvert^2=1.
\tag{8.6.14}\label{eq:8.6.14}
\end{equation}
To calculate the $S$-matrix element for absorbing or emitting such a photon,
we simply replace $e_\mu(\mathbf p,\mathord\pm1)$ in the Feynman rules with
\begin{equation}
e_\mu(\mathbf p)=\alpha_+e_\mu(\mathbf p,+1)
+\alpha_-e_\mu(\mathbf p,-1).
\tag{8.6.15}\label{eq:8.6.15}
\end{equation}
The polarization vectors for definite helicity satisfy the normalization
condition
\begin{equation}
e_\mu^*(\mathbf p,\lambda')e^\mu(\mathbf p,\lambda)
=\delta_{\lambda'\lambda}
\tag{8.6.16}\label{eq:8.6.16}
\end{equation}
and therefore in general
\begin{equation}
e_\mu^*(\mathbf p)e^\mu(\mathbf p)=1.
\tag{8.6.17}\label{eq:8.6.17}
\end{equation}
The two extreme cases are \emph{circular polarization}, for which
$\alpha_-=0$ or $\alpha_+=0$, and \emph{linear polarization}, for which
$\lvert\alpha_+\rvert=\lvert\alpha_-\rvert=1/\sqrt2$. For linear
polarization, by an adjustment of the overall phase of the state \eqref{eq:8.6.13}, we
can make $\alpha_+$ and $\alpha_-$ complex conjugates, so that they can be
expressed as
\begin{equation}
\alpha_\pm=\exp(\mp i\phi)/\sqrt2.
\tag{8.6.18}\label{eq:8.6.18}
\end{equation}
Then in the Feynman rules we should use a polarization vector
\begin{equation}
e_\mu(\mathbf p)=R(\widehat{\mathbf p})
\begin{bmatrix}
0\\[1pt]\cos\phi\\[1pt]\sin\phi\\[1pt]0
\end{bmatrix}.
\tag{8.6.19}\label{eq:8.6.19}
\end{equation}
That is, $\phi$ is the azimuthal angle of the photon polarization in the plane
perpendicular to $\mathbf p$. Note that the photon polarization vector here is
real, which is only possible for linear polarization. In between the extremes
of circular and linear polarization are the states of \emph{elliptic
polarization}, for which $\lvert\alpha_+\rvert$ and
$\lvert\alpha_-\rvert$ are non-zero and unequal.

More generally, an initial photon may be prepared in a statistical mixture of
spin states. In the most general case, an initial photon may have any number of
possible polarization vectors $e_\mu^{(r)}(\mathbf p)$, each with probability
$P_r$. The rate for absorbing such a photon in a given process will then be of
the form
\begin{equation}
\Gamma=\sum_rP_r\lvert e_\mu^{(r)}(\mathbf p)M^\mu\rvert^2
=M^{\mu *}M^\nu\rho_{\nu\mu},
\tag{8.6.20}\label{eq:8.6.20}
\end{equation}
where $\rho$ is the \emph{density matrix}
\begin{equation}
\rho_{\nu\mu}\equiv\sum_rP_r
e_\nu^{(r)}(\mathbf p)e_\mu^{(r)*}(\mathbf p).
\tag{8.6.21}\label{eq:8.6.21}
\end{equation}
Since $\rho$ is obviously a Hermitian positive matrix of unit trace (because
$\sum_rP_r=1$) with $\rho_{\mu0}=\rho_{0\mu}=0$ and
$\rho_{\nu\mu}p^\nu=\rho_{\nu\mu}p^\mu=0$, it may be written as
\begin{equation}
\rho_{\nu\mu}=\sum_{s=1,2}\lambda_s
e_\nu(\mathbf p;s)e_\mu^*(\mathbf p;s),
\tag{8.6.22}\label{eq:8.6.22}
\end{equation}
where $e_\mu(\mathbf p;s)$ are the two orthonormal eigenvectors of $\rho$ with
\begin{equation}
e_0(\mathbf p;s)=e_\mu(\mathbf p;s)p^\mu=0
\tag{8.6.23}\label{eq:8.6.23}
\end{equation}
and $\lambda_s$ are the corresponding eigenvalues, with
\[
\lambda_s\geq0,\qquad \sum_{s=1,2}\lambda_s=1.
\]
We may then write the rate for the photon absorption process as
\begin{equation}
\Gamma=\sum_{s=1,2}\lambda_s
\lvert e_\nu(\mathbf p;s)M^\nu\rvert^2.
\tag{8.6.24}\label{eq:8.6.24}
\end{equation}
Thus any statistical mixture of initial photon states is always equivalent to
having just two orthonormal polarizations $e_\nu(\mathbf p;s)$ with
probabilities $\lambda_s$.

In particular, if we know nothing whatever about the initial photon
polarization, then the two probabilities $\lambda_s$ for the polarization
vectors $e_\nu(\mathbf p;s)$ are equal, so that
$\lambda_1=\lambda_2=\frac12$, and the density matrix (and hence the
absorption rate) is an average over initial polarizations
\begin{equation}
\rho_{ij}=\frac12\sum_{s=1,2}
e_i(\mathbf p;s)e_j^*(\mathbf p;s)
=\frac12(\delta_{ij}-\widehat p_i\widehat p_j).
\tag{8.6.25}\label{eq:8.6.25}
\end{equation}
Fortunately, this result does not depend on the particular pair of
polarization vectors $e_i(\mathbf p;s)$ over which we average; for unpolarized
photons we can average the absorption rate over any pair of orthonormal
polarization vectors. Similarly, if we make no attempt to measure the
polarization of a photon in the final state, then the rate may be calculated
by summing over any pair of orthonormal final photon polarization vectors.

The same remarks apply to electrons and positrons; if (as is usually the case)
we make no attempt to prepare an electron or positron so that some spin states
are more likely than others, then the rate is to be calculated by
\emph{averaging} over any two orthonormal initial spin states, such as those
with spin $z$-component $\sigma=\mathord\pm\frac12$; if we make no attempt to
measure a final electron's or positron's spin state, then we must \emph{sum}
the rate over any two orthonormal initial spin states, such as those with spin
$z$-component $\sigma=\mathord\pm\frac12$. Such sums may be performed using
the relations \eqref{eq:5.5.37} and \eqref{eq:5.5.38}:
\begin{equation}
\sum_\sigma u_\alpha(\mathbf p,\sigma)\overline u_\beta(\mathbf p,\sigma)
=\left(\frac{-i\sl{p}+m}{2p^0}\right)_{\alpha\beta},
\tag{8.6.26}\label{eq:8.6.26}
\end{equation}
\begin{equation}
\sum_\sigma v_\alpha(\mathbf p,\sigma)\overline v_\beta(\mathbf p,\sigma)
=\left(\frac{-i\sl{p}-m}{2p^0}\right)_{\alpha\beta},
\tag{8.6.27}\label{eq:8.6.27}
\end{equation}
where $p^0=\sqrt{\mathbf p^2+m^2}$. For instance, if the initial state
contains an electron with momentum $\mathbf p$ and spin $z$-component
$\sigma$, and a positron with momentum $\mathbf p'$ and spin $z$-component
$\sigma'$, then the $S$-matrix element for the process will be of the form
$\overline v_\alpha(\mathbf p',\sigma')\mathcal M_{\alpha\beta}
u_\beta(\mathbf p,\sigma)$. Hence if neither electron nor positron spins are
observed, the rate will be proportional to
\begin{align*}
&\frac14\sum_{\sigma',\sigma}
\left\lvert\overline v_\alpha(\mathbf p',\sigma')
\mathcal M_{\alpha\beta}u_\beta(\mathbf p,\sigma)\right\rvert^2
=
\frac14\operatorname{Tr}\left\{
\beta\mathcal M^\dagger\beta
\left(\frac{-i\sl{p'}-m}{2p'^0}\right)
\mathcal M
\left(\frac{-i\sl{p}+m}{2p^0}\right)
\right\}.
\end{align*}
Techniques for the calculation of such traces are described in the Appendix to
this chapter.
